\begin{textsample}{POS Dim 3 – rj – Score 13.00 – rj/rj\_inf\_intro\_3.txt}  \label{ex:f3_pos_016}
Educação \textbf{Infantil} é um Direito O lugar da \textbf{criança} brasileira na política pública de educação é o de sujeito histórico, protagonista e cidadão com direito à educação a partir do nascimento, em estabelecimentos educacionais instituídos com a função de cuidar e educar como um único e indissociável ato promotor de seu desenvolvimento integral, de forma global e harmônica, nos aspectos físico, social, afetivo e cognitivo.
A educação \textbf{infantil} é a primeira etapa da educação básica a que todo cidadão brasileiro tem direito e que o Estado tem obrigação de garantir sem exceção nem discriminação ( NUNES, CORSINO e DIDONET, 2011, p. 9 ).
Hoje a Educação \textbf{Infantil} é um direito da \textbf{criança} e continua sendo um direito da mãe trabalhadora.
Durante muito tempo, na sociedade brasileira, a proteção à \textbf{criança} e à mãe foi considerada uma benesse.
A benevolência de alguns médicos associada à antiga ideia da filantropia, exercida por religiosos e pessoas de boa vontade, sempre como iniciativas particulares, foram as ações desenvolvidas no sentido da proteção materno-infantil.
Nem sempre houve no Brasil uma legislação específica sobre a infância e a educação das \textbf{crianças} \textbf{pequenas}, o que permite dizer que, por muitos séculos, não se reconheceu o acesso às instituições de educação destinada a \textbf{pré-escolares} como um direito, seja ele da mãe ou da \textbf{criança} \textbf{pequena}, e sim como um favor.
Nas últimas décadas do século XX, a Educação \textbf{Infantil} ( \textbf{creches} e \textbf{pré-escolas} ) no Brasil teve avanços significativos na sua legislação, diretrizes e normas, como uma conquista dos movimentos sociais.
Também contribuíram para isso as novas concepções de infância e \textbf{criança} produzidas nos estudos e pesquisas da sociologia da infância, da antropologia, da filosofia, da história, da pedagogia da infância, entre outras.
A Educação \textbf{Infantil}, primeira etapa da Educação Básica, passou a ser considerada um direito das \textbf{crianças} e um dever do Estado e da sociedade, sendo integrada aos sistemas municipais de ensino com a finalidade de promover o desenvolvimento integral da \textbf{criança} de até cinco anos.
Ela é oferecida em \textbf{creches} e \textbf{pré-escolas}.
Estas se caracterizam como espaços institucionais não \textbf{domésticos}, ou seja, são estabelecimentos educacionais públicos ou privados que educam e cuidam de \textbf{crianças} de 0 a 5 anos de idade no período diurno, em jornada integral ou parcial, regulados e supervisionados por órgão competente do sistema de ensino e submetidos a controle social. ( BRASIL, 2009 ) A Constituição Federal de 1988 e a Lei de Diretrizes e Bases da Educação Nacional ( LDB ) determinam : a Educação \textbf{Infantil} é um direito das \textbf{crianças} ; embora não seja obrigatória de 0 a 3 anos, hoje a \textbf{pré-escola} é obrigatória para as \textbf{crianças} de 4 a 5 anos ( Lei no 12.796/2013 ).
De acordo com as Diretrizes Curriculares Nacionais para a Educação \textbf{Infantil} ( DCNEI ), no art. 5o [... ] § 2{\EmojiFont °} : “ É obrigatória a matrícula na Educação \textbf{Infantil} de \textbf{crianças} que completam 4 ou 5 anos até o dia 31 de março do ano em que ocorrer a matrícula ”.
E no § 3o : “ As \textbf{crianças} que completam seis anos após o dia 31 de março devem ser matriculadas na Educação \textbf{Infantil} ”.
Os sistemas de ensino têm autonomia para complementar a legislação nacional por meio de normas próprias, específicas e adequadas às características locais.
Respeitadas as diretrizes curriculares nacionais e as normas estabelecidas pelas secretarias de educação, os estabelecimentos de ensino têm autonomia pedagógica.
Cabe, portanto, às instituições de Educação \textbf{Infantil} elaborar e executar sua proposta pedagógica, considerando as aprendizagens essenciais : Cuidar e educar ampliando as ações de conhecer o mundo ; formações de vínculo ; incentivo à autonomia ; respeito ao tempo de cada \textbf{criança} ; planejamento de espaços para livre exploração, organização do tempo e \textbf{escuta} ativa.
Como consta nas Diretrizes Curriculares Nacionais da Educação Básica ( Brasil, 2013 ), dizendo de outro modo, nessa etapa deve -se assumir o \textbf{cuidado} e a educação, valorizando a aprendizagem para a conquista da cultura da vida, por meio de atividades \textbf{lúdicas} em situações de aprendizagem ( jogos e \textbf{brinquedos} ), formulando proposta pedagógica que considere o currículo como conjunto de experiências em que se articulam saberes da experiência e socialização do conhecimento em seu dinamismo, depositando ênfase : I – na gestão das emoções ; II – no desenvolvimento de \textbf{hábitos} higiênicos e alimentares ; III – na vivência de situações destinadas à organização dos objetos pessoais e escolares ; IV – na vivência de situações de preservação dos recursos da natureza ; V – no contato com diferentes linguagens representadas, predominantemente, por ícones – e não apenas pelo desenvolvimento da prontidão para a leitura e escrita –, como potencialidades indispensáveis à formação do interlocutor cultural.
No seu artigo 8o, a Resolução CEB/CNE no 05/09 dispõe sobre os objetivos da proposta pedagógica das instituições de Educação \textbf{Infantil} : [... ] deve ter como objetivo garantir à \textbf{criança} acesso a processos de apropriação, renovação e articulação de conhecimentos e aprendizagens de diferentes linguagens, assim como o direito à proteção, à saúde, à liberdade, à confiança, ao respeito, à dignidade, à \textbf{brincadeira}, à convivência e à interação com outras \textbf{crianças} ( BRASIL, 2009 ).

% matched lemmas: brincadeira, brinquedo, creche, criança, cuidado, doméstico, escuta, hábito, infantil, lúdico, pequeno, pré-escola, pré-escolar
\end{textsample}
